\documentclass[11pt]{article}
\usepackage[a4paper,margin=25mm]{geometry}
\usepackage{amsmath,amssymb,amsthm,booktabs,array}
\usepackage[colorlinks=true,urlcolor=blue,linkcolor=blue]{hyperref}
\newcommand{\Q}{\mathbb Q}\newcommand{\Z}{\mathbb Z}
\newcommand{\F}{\mathbb F}\newcommand{\OO}{\mathcal O}
\newcommand{\Sel}{\operatorname{Sel}}\newcommand{\Cl}{\operatorname{Cl}}
\newcommand{\rank}{\operatorname{rank}}\newcommand{\Jac}{\operatorname{Jac}}
\newcommand{\ns}{\mathrm{ns}}
\newtheorem{theorem}{Theorem}
\newtheorem{proposition}[theorem]{Proposition}
\newtheorem{lemma}[theorem]{Lemma}
\emergencystretch=3em
\title{Rational points on \(X_{\ns}^{+}(121)\)\\
\large Proof draft and computational certificates}
\author{}\date{30 October 2026}
\begin{document}\maketitle

\noindent
This note combines the four previously completed even cases with the
three odd cases. A new six-prime sieve replaces the inherited large
tables. The class-group input for the odd cases is proved using
unconditional discriminant bounds. The odd local calculations were run
at two precisions, and their finite linear algebra was audited separately.
The proof concerning formal groups is given below; it is a mathematical
argument, not a machine formalization.

\begin{theorem}\label{thm:main}
The curve \(X_{\ns}^{+}(121)\) has precisely seven rational points.
They are CM points, with CM discriminants
\[
 -3,\ -4,\ -12,\ -16,\ -27,\ -67,\ -163.
\]
There are no rational cusps and no non-CM rational points.
\end{theorem}

We use the following model and generator:
\[
 E=X_{\ns}^{+}(11):\quad y^2+y=x^3-x^2-7x+10,
 \qquad E(\Q)=\Z G,\quad G=(4,5).
\]
The seven points in the theorem have the following images under
\(X_{\ns}^{+}(121)\longrightarrow E\).
\begin{center}\begin{tabular}{rlll}\toprule
\(n\)&\(nG\)&CM discriminant&\(j\)\\\midrule
\(-2\)&\((2,-1)\)&\(-27\)&\(-12288000\)\\
\(-1\)&\((4,-6)\)&\(-16\)&\(287496\)\\
\(0\)&\(O\)&\(-12\)&\(54000\)\\
\(1\)&\((4,5)\)&\(-67\)&\(-147197952000\)\\
\(2\)&\((2,0)\)&\(-4\)&\(1728\)\\
\(3\)&\((5/4,7/8)\)&\(-3\)&\(0\)\\
\(4\)&\((-2,3)\)&\(-163\)&\(-262537412640768000\)\\\bottomrule
\end{tabular}\end{center}
The identification of the elliptic curve, its Mordell--Weil group and
the integral \(j\)-points agree with \cite{ST}. The supplied
\path{outputs/verify_model.m} checks the model and the displayed values.

\section{A six-prime sieve and the auxiliary curves}
Put \(K=\Q(z)\), where
\[
 z^5+z^4-4z^3-3z^2+3z+1=0,
 \qquad u=2-z.
\]
The unique prime above 11 is generated by \(u\). The five cusps are
\(r_i s\), with
\[
 (r_i)=(1,4,3,9,5),\qquad
 s=(-z^4-z^3+2z^2+3z+3,-2z^4-4z^3+6z^2+11z+1).
\]
Here \(s\) has order 11. There are no rational cusps, since this cusp
divisor is a closed point of degree five.

Suppose a rational point at level 121 maps to \(P=nG\). At any prime
away from 11 at which \(j(P)\) is nonintegral, its valuation is divisible
by 121. Indeed, after a quadratic twist the associated elliptic curve
is a Tate curve. A nontrivial unipotent inertia element modulo 121
cannot belong to the normalizer of an unramified Cartan: this holds
modulo 11, and an element \(1+11N\), with \(N\) nonzero nilpotent
modulo 11, cannot belong to the Cartan congruence kernel either.
The Tate parameter therefore has valuation divisible by 121.
As the polar multiplicity of \(j\) at each cusp of \(E\) is 11, the
intersection multiplicities of \(P\) with the cusps are divisible by 11
away from 11.

For \(m=1,2,3\), let \(a_{m,i}\) be the representative in
\(\{1,\ldots,10\}\) of \(r_i^m\bmod11\), and choose \(f_m\) with
\[
 \operatorname{div}(f_m)=\sum_i a_{m,i}(r_i s)-22(O),
 \qquad \left.f_m(x/y)^{22}\right|_O=1.
\]
The class of \(f_m(P)\) lies in \(K(S,11)\), where
\(S=\{(u)\}\). The class number of \(K\) is one. A basis for this
five-dimensional space is
\[
 u,\quad z^4-3z^2+1,\quad z^2+z-1,\quad z,
 \quad -z^4-z^3+3z^2+3z.
\]
At \(O\) one uses the normalized leading coefficient. This also explains
the zero entry in the local tables whenever \(nG\) reduces to \(O\).

The sieve starts with all three classes unknown. Local residue-symbol
equations, with cusp residues retained conservatively, give:
\begin{center}\begin{tabular}{rrrrl}\toprule
prime&order of \(\bar G\)&modulus&states&ranks of each class system\\\midrule
353&11&11&7&3\\
23&33&33&10&3 or 5\\
3301&33&33&7&5\\
12101&1089&1089&12&5\\
15731&726&2178&7&5\\
53&12&4356&7&5\\\bottomrule
\end{tabular}\end{center}
The final states are exactly
\begin{equation}\label{eq:sieve}
 n\equiv n_0\pmod{4356},\qquad -2\le n_0\le4,
 \qquad [f_m(nG)]=[f_m(n_0G)]\quad(m=1,2,3).
\end{equation}
All six local tables were regenerated from the functions above.
The program \path{six_prime_sieve.py} performs the exhaustive affine
linear algebra and congruence lifting using only standard-library Python.
The final classes agree with the exact number-field calculation in
\path{outputs/cm_classes.m}. No higher-prime table or height bound is used.

Set
\[
 \alpha=x(2s)=z^4-z^3-3z^2+3z+2,\qquad
 \delta_{n_0}=x(n_0G-3s)-\alpha.
\]
There is an exact normalized function identity
\[
 \frac{x(P-3s)-\alpha}{x(3s)-\alpha}
   = f_1(P)^7 f_2(P)^{10} f_3(P)^6 f_T(P)^3 h(P)^{11},
\]
where \(\operatorname{div}(f_T)=11(s)-11(O)\). This is verified in
\path{outputs/verify_descent.m}. The class of \(f_T(P)\) is an
unramified descent homomorphism on \(E(K)/11E(K)\).
Consequently \eqref{eq:sieve} implies the necessary equation
\begin{equation}\label{eq:cusp}
 x(nG-3s)=\alpha+\delta_{n_0}w^{11},\qquad w\in K.
\end{equation}

Let
\[
 L=K(\theta),\quad \theta^3-\theta^2-7\theta+41/4=0,
 \qquad \beta^2=\alpha-\theta,\quad \beta\in L,
 \qquad F=L(\zeta_{11}).
\]
The degrees of \(L,F\) are 15 and 30. The two-descent map \(x-\theta\)
on \(E(L)\) shows that \eqref{eq:cusp} gives a point on
\begin{equation}\label{eq:genus5}
 \eta v^2=\alpha-\theta+\delta_{n_0}w^{11},\qquad
 \eta=\begin{cases}1&n_0\text{ even},\\4-\theta&n_0\text{ odd}.
 \end{cases}
\end{equation}
Indeed, \(-3s\) is twice a point and \(n\equiv n_0\bmod2\).
Put \(a=\delta_{n_0}/\beta^2\) and \(b=a^5\). The ideal of \(a\)
in \(F\) is the square of its unique prime over 2. In particular it
is a unit at 11. Over \(B=F(\sqrt\eta)\), the substitutions
\[
 X=aw,\qquad Y=\sqrt\eta\,a^5v/\beta
\]
give the unit model
\begin{equation}\label{eq:unit}
 C:\quad Y^2=X^{11}+b^2.
\end{equation}
Its known point \(P_0\) has \(X=a\). We write \(J\) for its Jacobian
in an even case, and \(J^\eta\) for the quadratic twist over \(F\) in
an odd case. All these identities and field constructions are checked in
\path{outputs/retry/verify_genus5.m}.

\section{The odd global descent, without GRH}
In this section \(n_0=-1,1,3\) and \(\eta=4-\theta\). Let
\(\tau\) generate \(\operatorname{Gal}(B/F)\). The twist identifies
\(J^\eta(F)\) with the anti-invariant subgroup of \(J(B)\).
The endomorphisms
\[
 \sigma(X,Y)=(\zeta_{11}X,Y),\qquad \lambda=1-\sigma
\]
commute with twisting. The kernel of \(\lambda\) on \(J\) is generated
by \(T=[(0,b)-\infty]\), which has order 11. The action of \(\tau\)
on \(T\) is trivial, so \(J^\eta[\lambda](F)=0\).
The Kummer map is represented over \(B\) by \(Y-b\); its values belong
to the anti-invariant part of \(B^\times/B^{\times11}\).

The exact discriminants are
\[
 |\operatorname{disc}F|=2^{20}11^{27},\qquad
 |\operatorname{disc}B|=2^{60}11^{56}.
\]
Both root discriminants are less than 19. The row \(b=9\) of
Odlyzko's unconditional Table 4 \cite{Od} gives
\[
 |D_M|>53.047^{r_1}20.710^{2r_2}e^{-24.001}.
\]
For any number field of degree \(d\ge330\), this implies
\[
 |D_M|^{1/d}>20.710\,e^{-24.001/330}
 >20.710\left(1-\frac{24.001}{330}\right)>19.
\]
An everywhere unramified cyclic extension of degree 11 of \(F\) or
\(B\) would have degree at least 330 and the same root discriminant.
Thus
\begin{equation}\label{eq:classes}
 \Cl(F)[11]=\Cl(B)[11]=0.
\end{equation}
The program \path{verify_class_obstruction.m} checks the discriminants
and the rational inequality. The tables' lower/upper rounding conventions
are stated explicitly in \cite{OdDescription}.

Let \(S\) consist of the primes over 2 and 11. Each of \(F\) and \(B\)
has four such primes. Their unit ranks are 14 and 29, and both contain
\(\mu_{11}\). It follows from \eqref{eq:classes} that
\[
 \dim F(S,11)=19,\qquad \dim B(S,11)=34.
\]
Norm is onto, since norm composed with restriction is multiplication by
2. Therefore the anti-invariant global Kummer space has dimension 15.

To construct this space, the program discovers candidates in the
degree-30 fields
\[
 L(\sqrt{4-\theta}),\qquad L(\sqrt{-11(4-\theta)}).
\]
For every candidate it checks the ideal valuations outside \(S\),
maps it into \(B\), and takes its quotient by its conjugate. Norm is
then exactly one. Fifteen independent joint local images certify a
basis of the entire anti-invariant space. Although GRH is used for
candidate discovery, membership and independence, together with the
proved dimension, remove that assumption from the answer.

There are three primes \(v\mid11\) of \(F\), all with
\(F_v\simeq\Q_{11}(\zeta_{11})\). At two of them, denoted \(v_1,v_2\),
\(B/F\) is ramified quadratic; at the third it is unramified quadratic.
It is nonsplit at all three, so \(J^\eta[\lambda](F_v)=0\).
The determinant of \(\lambda\) on the tangent space has valuation five
in \(F_v\), and hence
\[
 \dim_{\F_{11}} J^\eta(F_v)/\lambda J^\eta(F_v)=5.
\]
The local images are generated by actual anti-invariant divisors
\(6(P-\tau P)\), with \(P\in C(B_v)\), until this dimension is reached.
This allows divisors of degree two, which are needed here.

Let \(\mathcal S^+\) be the intersection of the certified global
15-dimensional space with these three local images. Omitting the
condition at 2 makes this a Selmer upper bound, which suffices. The
computed results are
\begin{equation}\label{eq:selmer}
 \dim\mathcal S^+=1\quad(n_0=-1,1,3),\qquad
 \mathcal S^+\longrightarrow
 \prod_{i=1}^2 B_{v_i}^\times/B_{v_i}^{\times11}
 \text{ is injective}.
\end{equation}
The known point has a nonzero class in each of these lines. Thus the
actual Selmer group and \(J^\eta(F)/\lambda J^\eta(F)\) both have
dimension one. Since the rational \(\lambda\)-torsion is zero and
\(\lambda\) has cyclotomic norm 11, this forces rank ten over \(F\).
The nontrivial automorphism of \(F/L\) conjugates the cyclotomic
action, so its fixed subspace has half this dimension. The corresponding
Jacobian therefore has rank exactly five over \(L\).
The rational-point argument below does not use this rank assertion.

\section{The formal group at a singular reduction}
We give the local argument explicitly because the odd known point
reduces to the cusp of \eqref{eq:unit}, not to a smooth point.
Fix either ramified place, put \(k=F_v\), \(l=B_v\), and
\(\pi=1-\zeta_{11}\). Normalize the valuation on \(l\) by
\(v_l(\pi)=2\); thus \(v_l(11)=20\).
One may use
\[
 \varpi=\frac{\sqrt\eta}{\pi^2},\qquad
 v_l(\varpi)=1,\qquad \tau\varpi=-\varpi.
\]
Both residue fields are \(\F_{11}\).

Use the proper hyperelliptic model of \eqref{eq:unit}. Its special
fibre is geometrically integral, with a single cusp and a smooth point
at infinity. The formal completion at the identity of its relative
Picard group is a smooth formal group of dimension five. One can also
construct its formal coordinates from sums of Abel maps at five smooth
sections with distinct residual \(X\)-coordinates. The cotangent lattice
has basis
\[
 \omega_j=\frac{X^{j-1}\,dX}{2Y},\qquad 1\le j\le5.
\]
Write \(H_m\) for formal points whose coordinates lie in
\(\varpi^m\OO_l\).

\begin{lemma}\label{lem:log}
The formal logarithm and exponential have integral coefficients in
these integral coordinates. Moreover,
\[
 \lambda:H_m\xrightarrow{\sim}H_{m+2}\qquad(m\ge1).
\]
In logarithm coordinates dual to the \(\omega_j\), the action of
\(\lambda\) is diagonal with entries \(1-\zeta_{11}^j\).
\end{lemma}
\begin{proof}
The automorphism \(\sigma\) is the identity on the model modulo the
ideal \((\pi)\), including its nonreduced base change to
\(\OO_l/(\pi)\). Hence every coefficient of the formal power series
for \(1-\sigma\) is divisible by \(\pi\). Its linear part divided by
\(\pi\) is invertible, since
\((1-\zeta_{11}^j)/\pi\equiv j\pmod\varpi\).
Dividing the series by \(\pi\) and applying the formal inverse
function theorem proves the asserted isomorphism.

The norm correspondence \(1+\sigma+\cdots+\sigma^{10}\) is zero on
the Jacobian, since the quotient curve has genus zero. In
\(\Z[\sigma]\), multiplication by 11 is therefore a unit times
\(\lambda^{10}\). Every coefficient of its formal series is divisible
by \(\pi^{10}\), and hence by 11.
The logarithm identity
\(\log([11]t)=11\log(t)\) now proves integrality by induction on total
degree. In degree \(d\ge2\), the unknown coefficient is multiplied
by \(11^d-11\), whose quotient by 11 is a unit; the already known
terms are all divisible by 11. The inverse series is integral as well.
Finally \(\sigma^*\omega_j=\zeta_{11}^j\omega_j\), which gives the
diagonal action.
\end{proof}

\begin{lemma}\label{lem:singulardisc}
Suppose \(P_0=(a,Y_0)\), \(v_l(Y_0)=1\), and \(\bar a=-1\).
If \(P=(a+t,Y(t))\) is on the branch through \(P_0\) and
\(d=v_l(t)\ge28\), then
\[
 [P-P_0]\in H_{d-1}\setminus H_d,
 \qquad
 \log[P-P_0]=\frac{t}{2Y_0}(1,a,a^2,a^3,a^4)
                  \pmod{\varpi^d}.
\]
\end{lemma}
\begin{proof}
Expand the five differentials in the parameter \(t\), on the analytic
branch through \(P_0\). The coefficient of \(t^h\) has valuation at
least \(-2h-1\): use
\(Y(t)^{-1}=Y_0^{-1}(1+(f(a+t)-f(a))/Y_0^2)^{-1/2}\), whose binomial
coefficients are 11-integral. Upon integration, a term with \(h\ge1\)
has valuation at least
\[
 (h+1)d-2h-1-20v_{11}(h+1)>d-1.
\]
Thus the five integrals lie in \(\varpi^{d-1}\OO_l\), and their
leading terms are as asserted. Lemma~\ref{lem:log} allows their
exponential in the formal Picard group. This analytic map and the Abel
map have identical differentials and agree at \(P_0\); they agree on
the disc. The first leading coordinate is nonzero, proving the exact
depth. This argument does not require \(P_0\) to be a section through
the smooth locus of the model.
\end{proof}

For anti-invariant points one further has
\begin{equation}\label{eq:H2}
 H_2^- = H_3^-,\qquad \mu_v(H_2^-)=0.
\end{equation}
Indeed, \(\tau\) acts trivially on \(H_2/H_3\), because it fixes the
residue field and sends \(\varpi\) to \(-\varpi\). On an
anti-invariant point it acts by \(-1\). Since 2 is invertible, such a
point is in \(H_3\). The unique formal \(\lambda\)-division in
Lemma~\ref{lem:log} is again anti-invariant. In particular, the local
Kummer class of a point in \(H_1^-\) depends only on its leading
tangent vector in \(H_1/H_2\).

\section{Depth and signs on the odd rational-image locus}
Suppose \(n\ne n_0\), with \(n\equiv n_0\bmod4356\), and write
\(k_0=v_{11}(n-n_0)\ge2\). Over the completion of \(K\), normalize
\(v_K(u)=1\). The integral minimal elliptic model is
\[
 x'=(x-4)/u^2,\quad y'=(y+1/2)/u^3,\qquad
 (y')^2=(x')^3+\frac{11}{u^2}(x')^2+
           \frac{33}{u^4}x'+\frac{121}{4u^6}.
\]
Its discriminant has valuation three. The formal parameter of \(2G\)
is \(-x'/y'=4u\). The elliptic formal logarithm is an isometry on
this formal neighbourhood, since \(1>v_K(11)/10=1/2\). Thus
\[
 v_K(t'((n-n_0)G))=1+5k_0.
\]
For each odd \(n_0\), put \(q_0=n_0G-3s\). Exact calculations give
\[
 v_K(x'(q_0))=v_K(y'(q_0))=1,\qquad v_K(\delta_{n_0})=2.
\]
The first term in the formal translation of \(x'\) is a unit times
\(2y'(q_0)t'\). Higher terms have larger valuation. Consequently
\[
 v_K(x(nG-3s)-x(q_0))=4+5k_0,
\]
and \eqref{eq:cusp} gives
\begin{equation}\label{eq:odddepth}
 v_K(w-1)=5k_0-3,\qquad
 v_{F_v}(X-a)=10k_0-6,\qquad
 v_l(X-a)=20k_0-12.
\end{equation}
Here \(\bar w=1\), and the term \(11(w-1)\) uniquely determines the
valuation of \(w^{11}-1\).
One also checks \(\bar a=-1\) and \(v_l(Y_0)=1\) at the ramified
places. These are the inputs to Lemma~\ref{lem:singulardisc}.

It remains to select a single global square root whose signs agree
with \(P_0\) at both places. Let
\[
 R=((n-n_0)/2)G,\qquad Q=2R,\qquad
 h_\theta(R)=\frac{(x(R)-\theta)^2-(3\theta^2-2\theta-7)}{2y(R)+1}.
\]
Then \(h_\theta(R)^2=x(Q)-\theta\). If
\(m=(y(Q)-y(q_0))/(x(Q)-x(q_0))\), set
\[
 \rho(R)=
 -\frac{m(\theta-x(q_0))+y(q_0)+1/2}
          {(x(q_0)-\theta)h_\theta(R)}.
\]
The exact function-field identities are
\[
 \rho(R)^2=\frac{x(q_0+2R)-\theta}{x(q_0)-\theta},
 \qquad \rho(O)=1.
\]
Choose \(v=v_0\rho(R)\), where \(v_0\) is the known square root in
\eqref{eq:genus5}. This gives the required global choice.
Indeed, \(4\mid n-n_0\), so \(R\) is formal. Its original elliptic
parameter has \(F_v\)-valuation at least 20, while
\(v_{F_v}(x(q_0)-\theta)\le6\).
The expansion of \(\rho\) at \(O\) therefore has \(\rho(R)-1\) of
positive valuation (in fact at least 14). It chooses the branch through
\(P_0\) simultaneously. The identities and valuations are checked by
\path{verify_odd_depth.m}.

Let \(D=[P-P_0]\in J^\eta(F)\). By
Lemma~\ref{lem:singulardisc}, its depth at each ramified place is
\(20k_0-13\). After
\begin{equation}\label{eq:divisions}
 r=10k_0-7\equiv3\pmod{10}
\end{equation}
formal divisions by \(\lambda\), its depth is one. Up to a nonzero
scalar at each place, its leading tangent vector is
\begin{equation}\label{eq:direction}
 \mathbf d=
 \left(1,-\frac1{2^3},\frac1{3^3},-\frac1{4^3},\frac1{5^3}\right)
 =(1,4,9,6,3)\in\F_{11}^5.
\end{equation}
This follows from the diagonal action in Lemma~\ref{lem:log} and
\(\bar a=-1\). The exponent three is essential: the odd elliptic
point has \(v_K(y')=1\), unlike the even cases.
As a numerical cross-check, \path{verify_local_depth_examples.m} verifies
these valuations and the direction for \(n=n_0+4356\) and
\(n=n_0+11\cdot4356\), at both ramified places in all three cases.

\section{Computing and excluding the odd directions}
Here is an explicit computation of the Kummer image of
\eqref{eq:direction}. Choose smooth points with residual
\(X\)-coordinates \(r_i=0,2,3,4,8\), and square roots \(y_i\in k\);
at zero take \(y_i=-b\). The tangent matrix of the five Abel maps is
\[
 M_{ij}=\frac{r_i^{j-1}}{2\bar y_i}.
\]
It is invertible. Solve \(\mathbf uM=\mathbf d\) over \(\F_{11}\)
and lift the \(u_i\) to integers. Let
\[
 P_i^\pm=(r_i\pm\varpi u_i,y_i^\pm),
 \qquad y_i^\pm\equiv y_i\pmod\varpi.
\]
The anti-invariant formal divisor
\(6\sum_i(P_i^+-P_i^-)\) has leading tangent vector \(\mathbf d\),
because \(12\equiv1\bmod11\). Its Kummer image is therefore
\begin{equation}\label{eq:muformula}
 d_v=6\sum_i\bigl(\mu_v(y_i^+-b)-\mu_v(y_i^--b)\bigr).
\end{equation}
By \eqref{eq:H2}, this is the image of every anti-invariant formal
point with that leading tangent vector. The calculations check
\(d_v\ne0\) and membership in the full local image.

Let \(\mathcal S_R^+\) be the restriction of \(\mathcal S^+\) to the
two ramified places, and let
\[
 \mathcal D=\langle(d_{v_1},0),(0,d_{v_2})\rangle.
\]
The exact finite-field certificates give
\begin{center}\begin{tabular}{rcccc}\toprule
\(n_0\)&\(\dim\mathcal S^+\)&localization kernel&\(\dim\mathcal D\)&
\(\dim(\mathcal S_R^+\cap\mathcal D)\)\\\midrule
\(-1\)&1&0&2&0\\
\(1\)&1&0&2&0\\
\(3\)&1&0&2&0\\\bottomrule
\end{tabular}\end{center}
The calculations were performed at precisions 80 and 160. Every local
Kummer evaluation checks that its relative precision is at least
\(2v_l(11)+2\), which exceeds the threshold for deciding an eleventh
power. The file \path{audit_matrices.m} reconstructs the Selmer upper
bound and verifies these intersections from the saved integer matrices,
independently of the number-field and local-point computations.

\begin{proposition}\label{prop:odd}
If \(n_0\in\{-1,1,3\}\), \(n\equiv n_0\bmod4356\), and
\eqref{eq:cusp} holds, then \(n=n_0\).
\end{proposition}
\begin{proof}
Suppose otherwise and form \(D\) as above. Its local Kummer images at
the two ramified places vanish, since its formal depth is at least
three. By \eqref{eq:selmer}, its global Kummer class is zero. Hence
there is a global \(\lambda\)-division.
At each ramified place the division is unique, because the twisted
\(\lambda\)-torsion is zero. It is therefore exactly the formal
division of Lemma~\ref{lem:log}.
Repeat this argument while the depth is at least three. After the
\(r\) divisions in \eqref{eq:divisions}, one obtains a global point
whose pair of local Kummer images is in \(\mathcal S_R^+\).
Equations \eqref{eq:direction}--\eqref{eq:muformula} show that this
pair is a nonzero element of \(\mathcal D\), contradicting the
computed zero intersection.
\end{proof}

\section{The four even classes and the conclusion}
For completeness we recall the certificate and division argument for
the even classes. Its detailed local calculation is in
\path{outputs/continuation/even_classes.pdf}. That argument uses only
\(4\mid n-n_0\), \(121\mid n-n_0\), and \eqref{eq:cusp}; thus the
new modulus 4356 supplies all its hypotheses.

At each of the three places above 11, use \(k=F_v\), with uniformizer
\(\pi=1-\zeta_{11}\). The known point on \eqref{eq:unit} has
\(\bar X=8\) and \(\bar Y\ne0\). Now
\(\lambda:H_m\simeq H_{m+1}\). The formal image \(\mu(H_1)\) is the
one-dimensional unramified Kummer line. Put
\[
 U=\prod_{v\mid11}\mu_v(H_1),\qquad
 T_{\mathrm{loc}}=\prod_{v\mid11}\langle\mu_v(T)\rangle,\qquad
 T_{\mathrm{glob}}=\langle(\mu_v(T))_v\rangle.
\]
The certified Selmer localization is injective, and its image
\(\mathcal S\) satisfies
\[
 U\cap T_{\mathrm{loc}}=0,\qquad
 \mathcal S\cap(U+T_{\mathrm{loc}})=T_{\mathrm{glob}}.
\]
The Selmer dimensions for \(n_0=-2,0,2,4\) are \(5,5,4,4\).
These calculations do not depend on GRH: their global basis was
certified by exact ideal membership and independent local images.
The new discriminant argument also proves the needed
\(\Cl(F)[11]=0\) directly.

For \(k_0=v_{11}(n-n_0)\ge2\), the corresponding divisor has depth
\(d=10k_0-8\). The signs are fixed simultaneously by the exact formula
\[
 v(R)=\frac{(x(R)-\theta)^2-(3\theta^2-2\theta-7)}{2y(R)+1},
 \qquad R=(n/2)G+4s,
\]
which is a square root of \(x(nG-3s)-\theta\).
After \(d-1\equiv1\bmod10\) formal divisions, the unramified Kummer
coordinate is a nonzero scalar times
\[
 \sum_{j=1}^5\ell_{j-1}\,8^{j-1}j^{-1}=5\quad\text{in }\F_{11},
 \qquad (\ell_0,\ldots,\ell_4)=(1,2,2,1,9).
\]
These are the coefficients of the Lagrange polynomial selecting the
zero section from residual coordinates \(0,2,3,4,8\).

Global division proceeds as in the odd case, with a torsion correction.
A global preimage differs from the chosen formal preimages by local
multiples of \(T\). The intersection identity above forces these
multiples to agree and forces the formal Kummer images to be zero.
Subtract the common global multiple of \(T\), and continue. The final
nonzero formal Kummer image is a contradiction. Hence the even cases
also have \(n=n_0\). This is the isogeny division-image argument
underlying Selmer-group Chabauty; compare \cite{Stoll}.

It follows from the six-prime sieve, Proposition~\ref{prop:odd}, and
the even argument that the image of every rational point at level 121
is one of the seven points in the first table.
Each has a unique rational lift. To see this, 11 is inert in each of
the seven class-number-one CM orders, so the CM algebra supplies a
Cartan structure at level 121. The lifts of the level-11 Cartan algebra
form a two-dimensional \(\F_{11}\)-space. The canonical CM lift is
zero, and the norm-one Cartan acts freely on the nonzero vectors,
through a cyclic group of order 12. Modulo elliptic automorphisms,
the nonzero orbit lengths are 12, 6, or 4. CM reciprocity supplies
this action, so only the canonical lift is rational. The finite orbit
calculation is checked in \path{outputs/verify_cm_lifts.m}.
This proves Theorem~\ref{thm:main}.

\section{Reproduction and scope of the certificates}
All commands are run from the directory containing \path{outputs/}.
The full instructions are in \path{outputs/odd/README.txt}.
The principal new calculations are:
\begin{verbatim}
magma -b outputs/odd/generate_six_prime_sieve.m \
  > outputs/odd/six_prime_sieve_data.txt
python3 outputs/odd/six_prime_sieve.py
magma -b outputs/odd/verify_class_obstruction.m
magma -b outputs/odd/verify_odd_depth.m
magma -b PADPREC:=160 OUTTAG:=rerun outputs/odd/odd_selmer.m \
  outputs/odd/formal_directions.m
\end{verbatim}
The exact finite matrix audits run separately on the saved Selmer
and direction data. The even-case reproduction instructions and
certificates are retained in \path{outputs/continuation/}.
No files in the original \path{/home/fnajman/121} project were changed.
The older reports are retained as records of the earlier partial work;
their statements that the odd cases or completeness remain unresolved
are superseded by this note.

\begin{thebibliography}{9}
\bibitem{ST} R. Schoof and N. Tzanakis,
\emph{Integral points of a modular curve of level 11},
Acta Arith. \textbf{152} (2012), 39--49.
\url{https://www.mat.uniroma2.it/~schoof/AA-schoof-tzanakis.pdf}.
\bibitem{Od} A. Odlyzko, \emph{Table 4: Unconditional bounds for
discriminants}, 29 November 1976, row \(b=9\).
\url{https://www-users.cse.umn.edu/~odlyzko/unpublished/discr.bound.table4}.
\bibitem{OdDescription} A. Odlyzko, \emph{Discriminant bounds:
description of tables}, 29 November 1976.
\url{https://www-users.cse.umn.edu/~odlyzko/unpublished/discr.bound.tables.txt}.
\bibitem{Stoll} M. Stoll, \emph{Chabauty without the Mordell--Weil group}.
\url{https://www.mathe2.uni-bayreuth.de/stoll/papers/ratpts-selmer-2016-12-07.pdf}.
\end{thebibliography}
\end{document}
